\subsection{Indirect Competitors}\label{subsect:IndirectCompetitors}
As defined in our theoretical framework, indirect competitors are applications that target a similar user group but offer a different value proposition, or satisfy only a partial need. For TrainHub, these are applications and tools that focus exclusively on a single aspect of the user's fitness journey (e.g., only diet or only workout tracking) and lack any integration with the services of a physical gym facility.

Our research highlighted a strong fragmentation in the digital habits of gym-goers. We categorized these indirect competitors into four main groups:

\begin{itemize}
    \item \textbf{Workout Tracking (Hevy, GymLife, Bodymatica):}
    these applications allow users to build custom workout routines and track their weightlifting progress.
    \textit{Why they are indirect:} they operate as excellent standalone diaries, but they do not offer a direct, integrated communication channel with the user's actual Personal Trainer, nor do they support facility-specific features like class bookings or gym access;

    \item \textbf{Nutrition Tracking (MyFitnessPal, Yazio):}
    these are market leaders in their specific niche, highly specialized in calorie counting and macronutrient tracking.
    \textit{Why they are indirect:} they focus exclusively on the dietary aspect of fitness, completely detached from the user's workout routines and the gym's ecosystem;

    \item \textbf{Cardio \& General Tracking (Garmin Connect, Strava, Polar Flow, Suunto):}
    these platforms are deeply tied to wearable devices (smartwatches) and excel at tracking outdoor activities, heart rate, and overall health metrics.
    \textit{Why they are indirect:} they are generic health trackers. While extremely useful, they are entirely disconnected from the specific concepts of "gym management", "PT appointments", or "facility interaction";

    \item \textbf{Generic \& Analog Tools (Apple Notes, Excel/Numbers, PDFs, WhatsApp/Telegram, Emails):}
    interestingly, our market research revealed that a significant portion of both clients and professionals still rely heavily on these non-specific tools to log progress, share workout PDFs, and communicate. 
    \textit{Why they are indirect:} they represent the current, inefficient "status quo." Using three or four different generic apps to manage one fitness journey creates friction.
\end{itemize}

\paragraph{Key Takeaway}
The analysis of these indirect competitors strongly validates the core idea behind TrainHub. While users have access to excellent specialized tools, the real frustration lies in the fragmentation. TrainHub aims to beat this "status quo" by gathering the best aspects of these scattered tools and unifying them into a single, automated platform directly connected to the user's physical gym.